\documentclass[journal]{IEEEtran}
\usepackage{amsmath,amssymb,bm,mathtools}

% ---- convenience macros ----
\newcommand{\R}{\mathbb{R}}
\newcommand{\Sym}{\mathbb{S}}
\newcommand{\Rp}{R_{\pi/2}}
\newcommand{\tr}{\operatorname{tr}}
\newcommand{\norm}[1]{\left\lVert #1 \right\rVert}

\begin{document}

\section{Estimator Structure and Robustness}

\subsection{Equivalence of the Symmetric-Formation Estimator and Local Least Squares}
\label{sec:equiv}

Let $\sigma\in C^{3}(\R^{2})$ be sampled by $m$ robots at positions
$r_{i}=c+\bar r_{i}$, with $s_{i}=\sigma(r_{i})$. Fit a local quadratic model
about $c$ in the basis
\begin{equation}
\phi(\bar r)=\big[\,1,\;x,\;y,\;\tfrac12 x^{2},\;xy,\;\tfrac12 y^{2}\,\big]^{\!\top},
\qquad \bar r=(x,y),
\end{equation}
chosen so that the coefficient vector is the local Taylor data itself,
$\theta=[\,\sigma(c),\;\nabla\sigma(c)^{\!\top},\;H_{11},\,H_{12},\,H_{22}\,]^{\!\top}$,
where $H=H_{\sigma}(c)$. With $\Phi=[\phi(\bar r_{1})\cdots\phi(\bar r_{m})]^{\!\top}$
and $s=[s_{1}\cdots s_{m}]^{\!\top}$, the ordinary least-squares (OLS) estimate is
$\hat\theta=(\Phi^{\!\top}\Phi)^{-1}\Phi^{\!\top}s$, and all structure resides in the
Gram (moment) matrix $G=\Phi^{\!\top}\Phi=\sum_{i}\phi(\bar r_{i})\phi(\bar r_{i})^{\!\top}$.

Consider the formation of [Bri\~n\'on-Arranz \emph{et al.}]: $N$ robots uniformly
spaced on a circle of radius $D$, $\bar r_{i}=D(\cos\varphi_{i},\sin\varphi_{i})$,
$\varphi_{i}=\varphi_{0}+2\pi i/N$, plus one robot at the center. The angular
identities $\sum_{i}e^{\,\jmath m\varphi_{i}}=0$ for $0<m<N$ annihilate every odd
moment and every anisotropic even moment up to order four once $N\ge5$, so $G$ is
block diagonal,
\begin{equation}
G=\operatorname{blkdiag}\!\Big(\,G_{\mathrm{tr}}\in\Sym^{3},\;\;
\tfrac{ND^{2}}{2}I_{2},\;\;\tfrac{ND^{4}}{8}\,\Big),
\end{equation}
on the index groups $\{1,\tfrac12x^{2},\tfrac12y^{2}\}$, $\{x,y\}$, and $\{xy\}$,
respectively. The OLS normal equations therefore \emph{decouple}, and the $6\times6$
solve reduces to scalar divisions:
\begin{align}
\widehat{\nabla\sigma}(c) &= \frac{2}{ND^{2}}\sum_{i=1}^{N} s_{i}\,\bar r_{i}, \\
\widehat{H}_{12} &= \frac{8}{ND^{4}}\sum_{i=1}^{N} s_{i}\,x_{i}y_{i},
\end{align}
while the $\{1,\tfrac12x^{2},\tfrac12y^{2}\}$ block, after eliminating $\sigma(c)$
using the \emph{center} sample $s_{c}$, returns the diagonal curvatures through
\begin{equation}
\label{eq:K}
3\widehat{H}+\Rp\,\widehat{H}\,\Rp^{\!\top}
=\frac{16}{ND^{4}}\sum_{i=1}^{N}\big(s_{i}-s_{c}\big)\,\bar r_{i}\bar r_{i}^{\!\top}
=:K .
\end{equation}
Because $H+\Rp H\Rp^{\!\top}=\tr(H)\,I$ for any symmetric $H$, the operator in
\eqref{eq:K} is the isotropic map
\begin{equation}
\mathcal{L}(H)\;=\;3H+\Rp H\Rp^{\!\top}\;=\;\tr(H)\,I+2H,
\end{equation}
which acts as a gain of $4$ on the trace (spin-$0$) component and $2$ on the
deviatoric (spin-$2$) component; it is thus invertible, with
\begin{equation}
\label{eq:Hinv}
\widehat{H}=\tfrac12\!\left(K-\tfrac14\tr(K)\,I\right).
\end{equation}
Equations \eqref{eq:K}--\eqref{eq:Hinv} are exactly the closed form of
[Bri\~n\'on-Arranz \emph{et al.}]. Hence the symmetric-formation estimator
\emph{is} the OLS estimator specialized to a configuration whose Gram matrix is
block diagonal; symmetry is what converts the matrix solve into averages. For
quadratic $\sigma$ the third-order remainder vanishes and $\widehat{H}=H$ exactly.

Two design facts follow directly. The threshold $N\ge5$ is the vanishing of the
$m=4$ (spin-$4$) harmonic, which removes the orientation dependence of the
deviatoric channel; $N\ge4$ suffices for the gradient (vanishing $m=3$).
Cocircularity ($\norm{\bar r_{i}}=D$ for all ring robots) makes the column
$\tfrac12x^{2}+\tfrac12y^{2}$ collinear with the constant column, so the trace
(Laplacian) becomes identifiable only after an off-circle sample---the center
robot---is added.

\subsection{Robustness to Formation Error and the Centralization Coupling}
\label{sec:robust}

Let the realized positions be $r_{i}=r_{i}^{\circ}+\delta_{i}$, with $r_{i}^{\circ}$
the nominal symmetric positions. The closed form \eqref{eq:K}--\eqref{eq:Hinv}
forms $K$ from the realized $\bar r_{i}$ but inverts the \emph{ideal} operator
$\mathcal{L}$. Under perturbation the moment cancellations of
Section~\ref{sec:equiv} no longer hold: the gradient term re-appears and the
operator is perturbed, $K=\mathcal{L}_{\delta}(H)+B(\nabla\sigma,\delta)+\Psi_H$,
so that
\begin{equation}
\widehat{H}=H+\mathcal{L}^{-1}\!\big(\Delta\mathcal{L}(\delta)\,H+B(\nabla\sigma,\delta)\big)+\mathcal{O}(MD),
\end{equation}
a \emph{bias} of order $\norm{\delta}/D$ that has nonzero expectation and is not
reduced by averaging over measurement-noise realizations. (Numerically, a
zero-mean position jitter of standard deviation $0.08D$ produced errors of
$\sim40\%$ in $H_{11}$ and $\sim65\%$ in $H_{12}$ on a pure quadratic field.)

If, instead, the realized positions are \emph{known}, OLS uses the true Gram
matrix $G(r)=\sum_{i}\phi(\bar r_{i})\phi(\bar r_{i})^{\!\top}$. For quadratic
$\sigma$ then $\hat\theta=\theta$ exactly for \emph{any} non-degenerate
configuration, and under i.i.d.\ measurement noise OLS is the best linear unbiased
estimator (Gauss--Markov). Formation error degrades only the conditioning of
$G(r)$ (curable by nondimensionalizing the basis columns by their natural powers
of $D$), never introducing bias.

Crucially, OLS does not require centralization. It depends on the data only
through the additive sufficient statistics
\begin{equation}
A=\sum_{i}\phi(\bar r_{i})\phi(\bar r_{i})^{\!\top}\in\Sym^{6},
\qquad b=\sum_{i}\phi(\bar r_{i})\,s_{i}\in\R^{6},
\end{equation}
each obtainable by average consensus over any connected communication graph with
$\mathcal{O}(1)$ message size ($27$ scalars), after which every node solves the
common $6\times6$ system locally. An unbiased estimate is thus available with no
central node, provided each robot can form $\phi(\bar r_{i})$, i.e.\ knows its
position in a shared frame.

This exposes a coupling between the two assumptions. The closed form's distinctive
advantage is that it uses the \emph{prescribed} role positions $r_{i}^{\circ}$
rather than measured ones; but this is valid only when $\norm{\delta_{i}}$ is
negligible, which itself requires relative localization and formation control of
sufficient quality---precisely what decentralized OLS needs to populate
$\phi(\bar r_{i})$. Consequently, an inability to centralize does not by itself
favor the closed form: decentralized OLS already covers that case. What favors the
closed form is the strictly narrower regime in which positions can be
\emph{commanded} accurately but not \emph{measured} relatively (e.g.\ timed
circular sweeps with known phase $\varphi_{i}(t)=\omega t$). Whenever localization
is good enough to certify the symmetry the closed form assumes, it is good enough
to run decentralized OLS, which is unbiased on the realized geometry.

\subsection{Extension to Vector Fields: the Jacobian and its Determinant}
\label{sec:jacobian}

Let $F=(u,v):\R^{2}\to\R^{2}$ and $J=\nabla F=\!\begin{bsmallmatrix}u_{x}&u_{y}\\v_{x}&v_{y}\end{bsmallmatrix}$
(in general nonsymmetric). Applying the scalar estimator of
Section~\ref{sec:equiv} to each component shares the same design matrix $\Phi$, so
the stacked system is block diagonal and inherits all of its properties. A full
quadratic per component has $6$ coefficients; two components give $12$ unknowns,
while each robot supplies the pair $(u_{i},v_{i})$, i.e.\ two equations. Hence
\emph{six robots in general position exactly determine the quadratic vector field}
(for quadratic $F$), yielding $J(c)$, the component Hessians (equivalently
$\nabla J$), and therefore
\begin{equation}
\det J(c)=u_{x}v_{y}-u_{y}v_{x},
\qquad
\nabla\!\det J(c)\ \text{(by the product rule on }\nabla J).
\end{equation}
The pentagon-plus-center attains this minimum isotropically and in a single
determined solve.

The lower-order needs are smaller. The Jacobian $J(c)$ alone is first-order: a
ring of $N\ge4$ recovers $\nabla u,\nabla v$ exactly via the gradient estimator
(no center robot), hence $J(c)$ and $\det J(c)$. The minimum for a single
determined plane fit of each component is $N=3$ (a triangle), which returns $J$ at
the triangle \emph{centroid} with an $\mathcal{O}(D)$ curvature bias, since the
plane-fit gradient is not the true gradient at the centroid for a curved field.

At an equilibrium $F(c^{\star})=0$, the sign of $\det J(c^{\star})$ classifies the
critical point ($\det J<0\Leftrightarrow$ saddle), while $\nabla\!\det J$ and the
eigenstructure of $J$ govern the motion of $c^{\star}$ and the local separatrices.
The six-robot pentagon-plus-center therefore supplies, in one determined solve,
the quantities required to locate and track a moving saddle.

Finally, an information bound bears on multi-point Jacobian estimation. $N$ robots
provide exactly $2N$ scalar samples of $F$. Any set of Jacobian estimates at
$k>1$ points obtained from \emph{overlapping} robot subsets are linear images of
these same $2N$ samples and are therefore statistically correlated; treating them
as independent---e.g.\ refitting over triple-centroids that share robots---
understates their covariance and, for non-quadratic $F$, injects a
model-dependent bias. Independent Jacobian estimates at $k$ points require
disjoint subsets, hence $\ge 3k$ robots. The two-stage
``plane-fit-then-fit-centroids'' pipeline is thus dominated by a single quadratic
fit over the robots, which returns $J$ and $\nabla J$ at any interior point with a
consistent covariance and no centroid bookkeeping.

\subsection{Closed-Form Curvature of the Determinant Field}
\label{sec:detcurv}

The second-order structure of $\det J$ need not be estimated by a separate
formation: it is fixed algebraically by the component fits of
Section~\ref{sec:jacobian}. We make this precise and then delimit its domain of
validity, since the distinction governs the required number of robots.

Write the fitted components in the monomial basis
$\{1,x,y,x^{2},xy,y^{2},x^{3},x^{2}y,xy^{2},y^{3}\}$ as $u=\sum_{k}a_{k}\phi_{k}$
and $v=\sum_{k}b_{k}\phi_{k}$, so that $(a_{3},a_{4},a_{5})$ and
$(b_{3},b_{4},b_{5})$ are the curvature (quadratic) coefficients and
$(a_{6},\dots,a_{9})$, $(b_{6},\dots,b_{9})$ the cubic ones.

\paragraph{Quadratic fit (six robots)}
When the components are fit at degree two, every first derivative
$u_{x},u_{y},v_{x},v_{y}$ is affine in $\bar r$; hence
$\det J=u_{x}v_{y}-u_{y}v_{x}$, a product of affine functions, is exactly
quadratic and its Hessian $H_{\det J}$ is \emph{constant} over the patch.
Differentiating twice gives, in basis-free form,
\begin{equation}
\label{eq:detHquad}
H_{\det J}=
\begin{pmatrix}
2\!\left(u_{xx}v_{xy}-u_{xy}v_{xx}\right) & u_{xx}v_{yy}-u_{yy}v_{xx}\\[3pt]
u_{xx}v_{yy}-u_{yy}v_{xx} & 2\!\left(u_{xy}v_{yy}-u_{yy}v_{xy}\right)
\end{pmatrix},
\end{equation}
or equivalently, in the fit coefficients,
\begin{equation}
\label{eq:detHminor}
H_{\det J}=4
\begin{pmatrix}
\begin{vmatrix}a_{3}&a_{4}\\ b_{3}&b_{4}\end{vmatrix} &
\begin{vmatrix}a_{3}&a_{5}\\ b_{3}&b_{5}\end{vmatrix}\\[10pt]
\begin{vmatrix}a_{3}&a_{5}\\ b_{3}&b_{5}\end{vmatrix} &
\begin{vmatrix}a_{4}&a_{5}\\ b_{4}&b_{5}\end{vmatrix}
\end{pmatrix},
\end{equation}
the three entries being the maximal minors of the curvature-coefficient matrix
$\left[\begin{smallmatrix}a_{3}&a_{4}&a_{5}\\ b_{3}&b_{4}&b_{5}\end{smallmatrix}\right]$.
Thus $H_{\det J}$ depends \emph{only} on the second-order coefficients of $u$ and
$v$; the gradient and constant terms cancel identically. The six-robot
pentagon-plus-center of Section~\ref{sec:jacobian} therefore yields $H_{\det J}$ in
closed form with no additional measurement.

\paragraph{Redundancy of a second stage}
Sampling $\det J$ at further interior points and refitting a quadratic recovers
\eqref{eq:detHminor} \emph{exactly}---$\det J$ is already quadratic---so the extra
stage contributes only round-off and, by the information bound of
Section~\ref{sec:jacobian}, statistically correlated samples: each sampled value is
a deterministic function of the same six measurements. The two-stage
``plane-fit-then-fit-centroids'' pipeline is strictly dominated by direct
differentiation of the component fit.

\paragraph{The third-order obstruction}
Identity \eqref{eq:detHquad} is exact \emph{for a quadratic field}, but a generic
$C^{3}$ field is modeled only to second order by a six-robot fit. Differentiating
$\det J$ twice \emph{without} truncation yields two groups of terms; for the
$xx$ entry,
\begin{equation}
\label{eq:obstruction}
\partial^{2}_{xx}\!\det J=
\underbrace{2\!\left(u_{xx}v_{xy}-u_{xy}v_{xx}\right)}_{\text{curvature}\,\times\,\text{curvature}}
+\underbrace{\left(u_{xxx}v_{y}-v_{xxx}u_{y}\right)+\left(u_{x}v_{xxy}-v_{x}u_{xxy}\right)}_{\text{gradient}\,\times\,\text{third order}},
\end{equation}
and analogously for the other entries. A degree-two component fit sets all third
derivatives to zero and so returns only the curvature$\times$curvature part; the
omitted gradient$\times$third-order contribution is an $O(MD)$ bias that does not
average out under measurement noise. Restoring it requires the third derivatives of
$u$ and $v$, i.e.\ a degree-\emph{three} component fit: with $10$ coefficients per
component ($20$ unknowns) and two measurements per robot, the minimal determined
formation is $10$ robots. The full Hessian then reads, entrywise,
\begin{align}
\label{eq:detHcubic}
(H_{\det J})_{xx}&=4\!\left(a_{3}b_{4}-a_{4}b_{3}\right)
+2\!\left(a_{1}b_{7}-3a_{2}b_{6}+3a_{6}b_{2}-a_{7}b_{1}\right),\\
(H_{\det J})_{xy}&=4\!\left(a_{3}b_{5}-a_{5}b_{3}\right)
+2\!\left(a_{1}b_{8}-a_{2}b_{7}+a_{7}b_{2}-a_{8}b_{1}\right),\notag\\
(H_{\det J})_{yy}&=4\!\left(a_{4}b_{5}-a_{5}b_{4}\right)
+2\!\left(3a_{1}b_{9}-a_{2}b_{8}+a_{8}b_{2}-3a_{9}b_{1}\right),\notag
\end{align}
where the first term of each line is the quadratic result \eqref{eq:detHminor} and
the second is the cubic correction. (For cubic components $\det J$ is quartic, so
its Hessian is no longer constant; \eqref{eq:detHcubic} is its value at the
formation center.) Expression \eqref{eq:detHcubic} is exact whenever $u$ and $v$
are cubic and is the leading-order estimate otherwise.

\paragraph{Interpretation}
$H_{\det J}$ is the curvature of the scalar field $\det J$, whose zero level set
locates the degeneracies of the flow---the loci where $J$ becomes singular.
Together with $\det J(c)$ and $\nabla\!\det J(c)$ it furnishes a second-order local
model of that set, i.e.\ how the singular locus, and with it a tracked saddle and
its separatrices, bends near the formation center. The resolution hierarchy is
therefore: \emph{six} robots recover $J$, $\det J$, $\nabla\!\det J$, and the
curvature$\times$curvature part of $H_{\det J}$; \emph{ten} robots add the
third-order terms of \eqref{eq:detHcubic} and resolve $H_{\det J}$ in full.

\subsection{Exactness of the Six-Robot Formation for Quadratic Fields}
\label{sec:exact_quadratic}

When the underlying vector field $F=(u,v)$ is strictly quadratic, the cubic and higher-order coefficients vanish identically. In this regime, the six-robot formation is not merely a second-order approximation, but is both necessary and sufficient to recover the curvature of the determinant field exactly.

To see this, consider the components of $F$ expressed in the purely quadratic basis $\{1, x, y, x^2, xy, y^2\}$. Each component is governed by six coefficients, yielding $12$ unknowns in total. Because each robot supplies two equations via the measured pair $(u_i, v_i)$, a configuration of six robots in general position---such as the pentagon-plus-center---forms a precisely determined system that extracts these coefficients with zero truncation error.

For a quadratic field, the first derivatives $u_x, u_y, v_x, v_y$ are strictly affine functions of the spatial coordinates. Consequently, the determinant $\det J = u_x v_y - u_y v_x$ is a product of affine functions, rendering it exactly quadratic. Its Hessian, $H_{\det J}$, is therefore a constant matrix over the entire domain.

Because $\det J$ contains no terms of degree higher than two, its second derivative depends exclusively on the curvature (quadratic) coefficients of $u$ and $v$; all gradient and constant terms cancel identically upon the second differentiation. Furthermore, because the true field possesses no cubic terms, the gradient-times-third-order obstruction identified in \eqref{eq:obstruction} evaluates to strictly zero. Therefore, the six-robot quadratic fit bypasses the need for a larger ten-robot formation, yielding the exact, unbiased $H_{\det J}$ in a single determined solve.

\end{document}
